\Needspace{4\baselineskip}
\begin{Example}
\label{ex:differentiation-operator-closure-sobolev}
We now return to the differentiation operator $(\mathcal{D}(D),D)$ defined in Example~\ref{ex:differentiation-operator-domain-and-unboundedness}. In our comments above, we alluded to the fact that the domain of the operator was not ``optimal,'' in that it did not include such obvious functions as piecewise differentiable functions. We can now see that the problem is that the operator is not closed. We construct the closure as we did in the proof of Lemma~\ref{lem:closable-operator-zero-sequence-criterion} above. That is, we define the domain $\mathcal{D}(\overline D)$ of the closure to be the functions $u\in L^2(0,L)$ such that there exists a sequence $u_n\in\mathcal{D}(D)$ and an element $v\in L^2(0,L)$ such that\begin{enumerate}\item $u_n\to u$ in $L^2(0,L)$ and\item $\dfrac{\mathrm{d}u_n}{\mathrm{d}x}\to v$ in $L^2(0,L)$.\end{enumerate}However, this is simply the definition of the Sobolev space\[\mathcal{D}(\overline D):=\{u\in H^1(0,L)\mid u(0)=0\},\]where the boundary condition is taken in the sense of trace.
\end{Example}

\Needspace{4\baselineskip}
\begin{Example}
\label{ex:laplacian-dirichlet-closure-partial-proof}
Closing the Laplacian operator is accomplished in much the same way. We consider the problem of Dirichlet conditions. Suppose there is a sequence $u_n$ in $\widetilde{\mathcal{D}}_D(\Delta)$ and functions $u$ and $v$ in $L^2(\Omega)$ such that
\begin{equation}\label{eq:laplacian-closure-function-convergence}
u_n\to u\text{ in }L^2(\Omega),
\end{equation}
and
\begin{equation}\label{eq:laplacian-closure-image-convergence}
\Delta u_n\to v\text{ in }L^2(\Omega).
\end{equation}
In Chapter~\ref{chap:linear-elliptic-equations} we shall show that the extended domain is given by $\mathcal{D}_D(\Delta)=H^2(\Omega)\cap H_0^1(\Omega)$. Although we cannot completely justify this assertion at this time, we can give parts of the proof. First, we note that if $\Omega$ is of class $C^1$, we have $\mathcal{D}_D(\Delta)\supset H^2(\Omega)\cap H_0^1(\Omega)$. To see this, note $\mathcal{D}_D(\Delta)$ is a subset of the completion of $\widetilde{\mathcal{D}}_D(\Delta)$ in the $H^2(\Omega)$ norm. Also note that by Theorem~\ref{thm:zero-boundary-sobolev-characterization}, $H^2(\Omega)\cap H_0^1(\Omega)$ is the set of all $u\in H^2(\Omega)$ such that $u=0$ on $\partial\Omega$ in the sense of trace. However, using the trace theorem again, we see that this set is also the completion of $\widetilde{\mathcal{D}}_D(\Delta)$ in the $H^2(\Omega)$ norm. Second, we note that $\mathcal{D}_D(\Delta)\subset H^1(\Omega)$. To see this, note that we can use Green\textquotesingle s identity to get
\begin{equation}\label{eq:laplacian-closure-green-energy-identity}
\int_\Omega|\nabla u_n|^2=-\int_\Omega u_n\Delta u_n.
\end{equation}
Thus, an application of Hölder\textquotesingle s inequality shows that the sequence $u_n\in\mathcal{D}_D(\Delta)$ satisfying 1 and 2 is bounded in $H^1(\Omega)$. Thus, by Theorem~\ref{thm:alaoglu-weak-star-subsequence} (and the fact that $H^1(\Omega)$ is reflexive), $u_n$ has a weakly convergent subsequence $u_n\rightharpoonup\overline u$ in $H^1(\Omega)$. However, by the uniqueness of weak limits we must have $u=\overline u$.
\end{Example}

\Needspace{4\baselineskip}
\begin{Example}
\label{ex:differential-operators-closable-wave-domain}
Note that all differential operators with smooth coefficients are closable. To see this, suppose that $(\mathcal{D}(L),L)$ is a differential operator of order $m$ with $\mathcal{D}(L)\subset L^2(\Omega)$. Then if $x_n\to0$ in $L^2(\Omega)$, we have $L(x_n)\to0$ in $H^{-m}(\Omega)$. Since $L^2(\Omega)$ is continuously imbedded into $H^{-m}(\Omega)$, either $\lim L(x_n)$ (thought of as a limit in $L^2(\Omega)$) is equal to $0$ or does not exist. Thus, by Lemma~\ref{lem:closable-operator-zero-sequence-criterion}, the operator is closable.

Although all differential operators are closable, the domain of the closure is not necessarily a Sobolev space as one might expect. For instance, consider the wave operator
\begin{equation}\label{eq:wave-operator-closable-nonsobolev-example}
L(u):=u_{tt}-u_{xx},
\end{equation}
with domain $H^2(\mathbb{R}\times(0,T))$, with $T<\infty$. The reader should verify that such obvious wave-like solutions as $u(x,t)=f(x-t)$, where $f\in L^2(\mathbb{R})$ is discontinuous, are in the closure of the operator. Of course, such functions are neither in $H^2(\mathbb{R}\times(0,T))$ nor $H^1(\mathbb{R}\times(0,T))$. (They do not have a well defined trace along the lines of discontinuity.)
\end{Example}

\Needspace{4\baselineskip}
\begin{Example}
\label{ex:nonclosable-integral-operator}
While most operators we encounter in practice are closable, there are examples of operators which are not. Let $X=L^2(\mathbb{R})$, $Y=\mathbb{R}$, and define $(\mathcal{D}(A),A)$ by letting $\mathcal{D}(A)$ be the set of bounded functions with compact support and $Au:=\int_{-\infty}^\infty u(x)\,\mathrm{d}x$. Then if we let
\begin{equation}\label{eq:nonclosable-integral-operator-sequence}
u_n(x)=\begin{cases}\dfrac{1}{n},&|x|\leq n,\\0,&|x|>n,\end{cases}
\end{equation}
we see that $u_n\to0$ in $L^2(\mathbb{R})$ but $Au_n=2$.
\end{Example}
